
\section{\textbf{\texttt{exists()}} ORM Method}

\begin{frame}[fragile]{\textbf{\texttt{exists()}} Method}
	\begin{itemize}
		\item The \textbf{\texttt{exists()}} method checks whether the records in a recordset still exist in the database.
		\item It is useful when a recordset may contain deleted or invalid IDs.
		\item The method returns a recordset containing only the records that still exist.
\begin{block}{Method Syntax}
	\begin{itemize}
		\item The standard syntax of the exists() method is:
\begin{lstlisting}
def exists(self):
	return super().exists()
\end{lstlisting}
		\item This method has two important parts:
		\begin{enumerate}
			\item \textbf{\texttt{exists}}
			\item \textbf{\texttt{self}}
		\end{enumerate}
	\end{itemize}
\end{block}
	\end{itemize}
\end{frame}

\begin{frame}[fragile]{Breaking Down the Method Signature}
	\begin{block}{}
		\begin{itemize}
			\item \textbf{\texttt{exists}}: The ORM method that filters a recordset to real DB rows.
			\item \textbf{\texttt{self}}: The recordset you want to validate.
\begin{lstlisting}
students = self.env['student.student'].browse([15, 16, 9999])
print(students)
\end{lstlisting}
			Output is:
\begin{lstlisting}
student.student(15, 16, 9999)
\end{lstlisting}
			\item \texttt{browse()} does not check the database immediately.
			\item ID \texttt{9999} may not exist, but it is still inside the recordset.
		\end{itemize}
	\end{block}
\end{frame}

\begin{frame}[fragile]{What Does \texttt{exists()} Return?}
	\begin{block}{}
		\begin{itemize}
			\item \texttt{exists()} returns a recordset of only the records that exist.
\begin{lstlisting}
students = self.env['student.student'].browse([15, 16, 9999])
real_students = students.exists()
print(real_students)
\end{lstlisting}
			Output is (example):
\begin{lstlisting}
student.student(15, 16)
\end{lstlisting}
			\item ID \texttt{9999} was removed because it is not in the database.
		\end{itemize}
	\end{block}
\end{frame}

\begin{frame}[fragile]{Checking a Single Record}
	\begin{itemize}
		\item You can also test one record in an \texttt{if} statement.
\begin{lstlisting}
student = self.env['student.student'].browse(15)
if student.exists():
	print("Student is still in the database")
else:
	print("Student was deleted")
\end{lstlisting}
		\item An empty recordset is falsy in boolean context.
		\item This pattern avoids crashing when a record was already unlinked.
	\end{itemize}
\end{frame}

\begin{frame}[fragile]{When Should You Use \texttt{exists()}?}
	\begin{itemize}
		\item After long-running operations where records may have been deleted.
		\item When you received IDs from the UI, RPC, or another model.
		\item Before writing or unlinking a recordset that might be stale.
\begin{lstlisting}
records = self.env['student.student'].browse(ids).exists()
records.write({'active': True})
\end{lstlisting}
		\item \texttt{exists()} is a safety method: it does not create or update data.
	\end{itemize}
\end{frame}

\begin{frame}[fragile]{Method Call}
\begin{block}{Method Call}
\begin{lstlisting}
real_students = students.exists()
\end{lstlisting}
	\begin{itemize}
		\item Here, \textbf{\texttt{.exists()}} is the \textbf{method call}.
		\item No dictionary argument is required.
		\item The return value is always a recordset (possibly empty).
	\end{itemize}
\end{block}
\end{frame}
